\section{数据划分、泄漏与分布偏移}

模型在一份数据上表现良好，为什么能说明它会处理尚未见过的数据？关键条件是，评价数据在模型和选择规则固定之后，仍像目标部署分布的一次新抽样；单纯的``数据量大''不足以保证这一点。数据划分的作用，是在历史数据内部人为保留一次近似未来的遭遇。

\subsection{训练、验证和测试承担不同角色}

训练集用于拟合参数；验证集参与选择模型、特征、阈值和超参数；测试集只在整个选择过程结束后估计最终流程的表现。三者的差别不在文件名，而在信息是否参与过决定。

设候选流程为 \(\mathcal A_1,\ldots,\mathcal A_m\)。若用验证集选择

\[
\widehat j=\argmin_j \widehat R_{\mathrm{val}}(\mathcal A_j),
\]

那么最小验证误差也带有选择乐观性：候选越多，越容易挑到一个偶然适合验证噪声的流程。独立测试集的作用，是在 \(\widehat j\) 固定后重新取得一次没有参与选择的风险估计。

测试集不是天然永久有效。看过测试误差后修改特征、重新选择模型或调整阈值，测试数据就已经成为验证信息。重复查看不会在代码中留下``泄漏''标记，却会逐步消耗测试集的独立性。

\subsection{切分单位必须匹配独立单位}

随机按行切分只在行近似独立且部署也是同分布随机抽样时合理。若同一患者、用户、设备或地点贡献多行，应按群组切分；若预测未来，应按时间切分；若模型将部署到新医院或新地区，最好让验证或测试按域留出。

错误切分会让训练集和测试集共享模型在部署时得不到的信息。例如同一患者的相邻影像分别落入两侧，模型可能识别患者或设备特征，而不是疾病规律。测试分数很高，却只证明它能识别已经见过的群组。

时间数据还有方向性。用未来样本估计标准化参数、填补过去缺失值，或随机打乱后预测较早时刻，都让未来信息流向过去。即使标签列从未直接进入特征，也已经违反了真实预测时的信息边界。

\subsection{预处理也是训练的一部分}

标准化、缺失值填补、PCA、特征选择和类别重采样都会从数据估计参数。正确的评价对象是完整训练流程，不能只检查最后一个预测器：

\[
\mathcal A:S_{\mathrm{train}}\longmapsto \widehat f.
\]

因此每次验证切分中，预处理必须只在该次训练折上拟合，再原样应用于验证折。先在全数据上选择高相关特征，再做交叉验证，会让验证标签影响特征集合；先用全数据计算均值和主成分，也会让验证输入影响表示。

这类泄漏危险之处在于程序完全可以正常运行，甚至每一步单独看都符合公式。错误发生在信息流，而不是某个算术表达式。

\subsection{交叉验证估计的是一条训练规则}

\(K\) 折交叉验证轮流用 \(K-1\) 折训练、剩余一折评价，再平均验证损失。它复用了有限数据，适合比较流程和选择超参数，但估计对象是``用约 \((K-1)n/K\) 个样本训练的规则''在当前切分机制下的表现。

若超参数也通过交叉验证选择，而又希望无偏评价整个选择流程，可以使用嵌套交叉验证：内层选择，外层评价。外层测试折不能参与内层预处理或超参数决定。

交叉验证也不能创造缺失的独立单位。时间序列随机 \(K\) 折仍可能穿越时间，群组数据逐行 \(K\) 折仍可能共享个体。折数只是计算设计，切分机制才决定它模拟哪个未来。

\subsection{三种分布偏移改变不同条件分布}

将联合分布分解为

\[
P(X,Y)=P(Y\mid X)P(X)
       =P(X\mid Y)P(Y).
\]

不同偏移对应不同部分发生变化：

\begin{itemize}
\tightlist
\item
  covariate shift：\(P(X)\) 改变，而 \(P(Y\mid X)\) 暂时假设不变；
\item
  label shift：\(P(Y)\) 改变，而 \(P(X\mid Y)\) 暂时假设不变；
\item
  concept shift：\(P(Y\mid X)\) 本身改变，相同输入不再对应相同结果机制。
\end{itemize}

这些名称表示关于哪些条件分布保持不变的假设，无法只看直方图就确认。重要性加权可以在支持覆盖和密度比可靠时修正某些 covariate shift，却不能修复新分布出现在训练支持之外，也不能自动处理条件机制改变。

部署还可能由模型自己改变。推荐系统决定用户看见什么，风控模型改变哪些交易被人工审核，强化学习策略更直接地决定下一批状态。此时数据不再是被动环境的固定抽样，评价必须考虑策略引起的反馈。

\subsection{测试分数也有抽样误差}

即使测试集完全独立，报告的经验风险仍是随机量。若逐样本损失方差有限，均值的标准误大致按 \(1/\sqrt n\) 缩小；分类错误率接近 \(p\) 时，独立 Bernoulli 近似给出标准误

\[
\sqrt{\frac{p(1-p)}{n}}.
\]

但样本相关、长尾损失、小群体样本稀少或多次比较都会使这个简单公式失真。比较两个模型时，在同一测试样本上分析逐样本损失差通常比把两个独立标准误相加更有效，因为难例对两者造成的共同波动可以抵消。

小数点后更多位数不等于更多证据。报告评价结果时，至少应同时说明样本量、切分单位、时间范围、选择过程和不确定性；若部署分布可能变化，还应说明这份估计在哪些变化下失效。

\subsection{评价纪律的边界}

严格划分训练、验证和测试可以防止同一份随机性被重复利用，却不能保证标签定义正确、样本覆盖目标人群或未来机制稳定。一个完全没有泄漏的实验仍可能精确回答错误问题。

因此可信评价有两层：统计层要求保留独立信息并量化抽样误差；建模层要求评价数据真正模拟部署时的对象、时间与行动。前者防止把偶然性当规律，后者防止把旧世界的规律误当成新世界的保证。
